% from ega2-1-fr.tex
\oldpage[II]{173}
On notera aussi qu'il résulte du diagramme commutatif
(\hyperref[II.8.5.3.2-fr]{8.5.3.2}) que \emph{la restriction à $C_X$ de
$\operatorname{Proj}(\widehat{\alpha})$ n'est autre que le morphisme
$\operatorname{Spec}(\alpha)$}.
\end{proof}

\begin{corollary}[8.6.3]
\label{II.8.6.3-fr}
Considérés comme $X$-schémas, $\widehat{E}_X$ est canoniquement isomorphe à
$\operatorname{Spec}(\mathcal{S}_X^{\leq})$, et $E_X$ à
$\operatorname{Spec}(\mathcal{S}_X)$.
\end{corollary}

\begin{proof}
Comme on sait que les morphismes $p_X$ et $\pi_X$ sont affines
(\hyperref[II.8.3.5-fr]{8.3.5} et
\hyperref[II.8.3.6-fr]{8.3.6}), il suffit (vu
(\hyperref[II.1.3.1-fr]{1.3.1})) de vérifier le corollaire lorsque
$Y=\operatorname{Spec}(A)$ est affine et $\mathcal{S}=\widetilde{S}$. La
première assertion résulte de l'existence de l'isomorphisme canonique
(\hyperref[II.8.2.7.2-fr]{8.2.7.2})
$(S_f^{\geq})_{f/1}^{\leq}\xrightarrow{\sim}S_f^{\leq}$ et du fait que ces
isomorphismes sont compatibles avec le passage de $f$ à $fg$ ($f$ et $g$
homogènes dans $S_+$). De même, la formule
(\hyperref[II.8.3.6.2-fr]{8.3.6.2}), appliquée à $\pi_X$ au lieu de $\pi$,
montre que
$\pi_X^{-1}(u^{-1}(X_f))=
\operatorname{Spec}((S_f^{\geq})_{f/1})$ pour $f$ homogène dans $S_+$, et la
seconde assertion découle de l'existence de l'isomorphisme canonique
(\hyperref[II.8.2.7.2-fr]{8.2.7.2})
$(S_f^{\geq})_{f/1}\xrightarrow{\sim}S_f$.

On peut ainsi dire que $\widehat{C}_X$, considéré comme $X$-schéma, s'obtient
par \emph{recollement} des deux $X$-schémas affines sur $X$,
$C_X=\operatorname{Spec}(\mathcal{S}_X^{\geq})$ et
$\widehat{E}_X=\operatorname{Spec}(\mathcal{S}_X^{\leq})$, dont
l'intersection est l'ouvert $E_X=\operatorname{Spec}(\mathcal{S}_X)$.
\end{proof}

\begin{corollary}[8.6.4]
\label{II.8.6.4-fr}
Supposons que $\mathcal{O}_X(1)$ soit un $\mathcal{O}_X$-Module inversible et
que $\mathcal{S}_X$ soit isomorphe à
$\bigoplus_{n\in\mathbf{Z}}(\mathcal{O}_X(1))^{\otimes n}$ (ce qui sera en
particulier le cas si $\mathcal{S}$ est engendrée par $\mathcal{S}_1$
(\hyperref[II.3.2.5-fr]{3.2.5} et
\hyperref[II.3.2.7-fr]{3.2.7})). Alors le cône projectif épointé
$\widehat{E}$ s'identifie au fibré vectoriel de rang un
$\mathbf{V}(\mathcal{O}_X(-1))$ sur $X$, et le cône affine épointé $E$ au
sous-préschéma de ce fibré vectoriel induit sur le complémentaire de la
section nulle. Avec cette identification, la rétraction canonique
$\widehat{E}\to X$ s'identifie au morphisme structural du $X$-schéma
$\mathbf{V}(\mathcal{O}_X(-1))$. Enfin, il existe un $Y$-morphisme canonique
$\mathbf{V}(\mathcal{O}_X(1))\to C$, dont la restriction au complémentaire de
la section nulle de $\mathbf{V}(\mathcal{O}_X(1))$ est un isomorphisme de ce
complémentaire sur le cône affine épointé $E$.
\end{corollary}

\begin{proof}
En effet, si on pose $\mathcal{L}=\mathcal{O}_X(1)$,
$\mathcal{S}_X^{\geq}$ est alors identique à
$\mathbf{S}_{\mathcal{O}_X}(\mathcal{L})$, donc $\widehat{E}_X$ s'identifie
canoniquement à $\mathbf{V}(\mathcal{L}^{-1})$ en vertu de
(\hyperref[II.8.6.3-fr]{8.6.3}) et $C_X$ à $\mathbf{V}(\mathcal{L})$. Le
morphisme $\mathbf{V}(\mathcal{L})\to C$ est la restriction de
$\operatorname{Proj}(\widehat{\alpha})$, et les assertions du corollaire sont
donc des cas particuliers de (\hyperref[II.8.6.2-fr]{8.6.2}).
\end{proof}

On notera que l'image réciproque par le morphisme
$\mathbf{V}(\mathcal{O}_X(1))\to C$ de l'espace sous-jacent au préschéma des
sommets de $C$ est l'espace sous-jacent à la section nulle de
$\mathbf{V}(\mathcal{O}_X(1))$ (\hyperref[II.8.5.5-fr]{8.5.5})~; mais en
général les sous-préschémas correspondants de $C$ et de
$\mathbf{V}(\mathcal{O}_X(1))$ ne sont pas isomorphes. Cette question va être
étudiée ci-dessous.

\subsection{Éclatement de cônes projetants}
\label{subsection:II.8.7-fr}

\begin{env}[8.7.1]
\label{II.8.7.1-fr}
Sous les conditions de (\hyperref[II.8.6.1-fr]{8.6.1}), on a, en posant
$r=\operatorname{Proj}(\widehat{\alpha})$, un diagramme commutatif
\begin{equation}
\label{II.8.7.1.1-fr}
  \xymatrix{
    X\ar[r]^{i_X\circ\varepsilon_X}\ar[d]_{q} &
    \widehat{C}_X\ar[d]^{r}\\
    Y\ar[r]_{i\circ\varepsilon} & \widehat{C}
  }
\tag{8.7.1.1}
\end{equation}

